% from ega2-1-fr.tex
\oldpage[II]{92}

  \item[\rm{(vi)}] Soient $f:X\to Y$ un morphisme quasi-compact, $j$
  l'injection canonique $X_{\mathrm{red}}\to X$. Si $\mathcal{L}$ est un
  $\mathcal{O}_X$-Module ample relativement à $f$, alors
  $j^*(\mathcal{L})$ est ample relativement à $f_{\mathrm{red}}$.
\end{enumerate}
\end{proposition}

Notons d'abord que (v) et (vi) se dérivent de (i), (i bis) et (iv) par le
même raisonnement que dans (\hyperref[II.4.4.10-fr]{4.4.10}), en utilisant
(\hyperref[II.4.6.4-fr]{4.6.4}) au lieu de
(\hyperref[II.4.4.5-fr]{4.4.5})~; nous laissons le détail du raisonnement au
lecteur. L'assertion (i) est trivialement conséquence de
(\hyperref[II.4.4.10-fr]{4.4.10, (i)}) et de
(\hyperref[II.4.6.2-fr]{4.6.2}). Pour démontrer (i bis), (iii) et (iv), nous
utiliserons le lemme suivant~:

\begin{lemma}[4.6.13.1]
\label{II.4.6.13.1-fr}
\begin{enumerate}
  \item[\rm{(i)}] Soient $u:Z\to S$ un morphisme, $\mathcal{L}$ un
  $\mathcal{O}_S$-Module inversible, $s$ une section de $\mathcal{L}$
  au-dessus de $S$, $s'$ la section de $u^*(\mathcal{L})=\mathcal{L}'$
  au-dessus de $Z$ qui lui correspond canoniquement. On a alors
  $Z_{s'}=u^{-1}(S_s)$.
  \item[\rm{(ii)}] Soient $Z$, $Z'$ deux $S$-préschémas, $p$, $p'$ les
  projections de $T=Z\times_S Z'$, $\mathcal{L}$ (resp.
  $\mathcal{L}'$) un $\mathcal{O}_Z$-Module (resp. un
  $\mathcal{O}_{Z'}$-Module) inversible, $t$ (resp. $t'$) une section de
  $\mathcal{L}$ (resp. $\mathcal{L}'$) au-dessus de $Z$ (resp. $Z'$), $s$
  (resp. $s'$) la section de $p^*(\mathcal{L})$ (resp.
  $p'^*(\mathcal{L}')$) au-dessus de $Z\times_S Z'$ qui lui correspond.
  Alors on a $T_{s\otimes s'}=Z_t\times_S Z'_{t'}$.
\end{enumerate}
\end{lemma}

Il résulte des définitions que l'on peut se ramener au cas où tous les
préschémas considérés sont affines. En outre, dans (i), on peut supposer
$\mathcal{L}=\mathcal{O}_S$~; l'assertion (i) résulte alors aussitôt de
(\hyperref[I.1.2.2.2-fr]{I, 1.2.2.2}). De même, dans (ii), on peut se
limiter au cas où $\mathcal{L}=\mathcal{O}_Z$,
$\mathcal{L}'=\mathcal{O}_{Z'}$, et alors l'assertion se réduit au lemme
(\hyperref[II.4.3.2.4-fr]{4.3.2.4}).

Démontrons alors (i bis). On peut supposer $Y$ affine
(\hyperref[II.4.6.4-fr]{4.6.4}), donc $\mathcal{L}$ ample
(\hyperref[II.4.6.6-fr]{4.6.6})~; lorsque $s$ parcourt l'ensemble réunion
des $\Gamma(X,\mathcal{L}^{\otimes n})$ ($n>0$), les $X_s$ forment une
base de la topologie de $X$ (\hyperref[II.4.5.2-fr]{4.5.2, a)}), donc par
hypothèse les $j^{-1}(X_s)$ forment une base de la topologie de $X'$~; on
conclut donc par le lemme
(\hyperref[II.4.6.13.1-fr]{4.6.13.1, (i)}) et par
(\hyperref[II.4.5.2-fr]{4.5.2, a)}) que $j^*(\mathcal{L})$ est ample.

Démontrons ensuite (iii). On peut de nouveau supposer $Y$ et $Y'$ affines
(\hyperref[II.4.6.4-fr]{4.6.4}), d'où il résulte que la projection
$h:X'\to X$ est affine (\hyperref[II.1.5.5-fr]{1.5.5}). Comme
$\mathcal{L}$ est ample (\hyperref[II.4.6.6-fr]{4.6.6}), lorsque $s$
parcourt l'ensemble des sections des $\mathcal{L}^{\otimes n}$ ($n>0$)
au-dessus de $X$ telles que $X_s$ soit affine, les $X_s$ recouvrent $X$
(\hyperref[II.4.5.2-fr]{4.5.2, a')})~; donc les $h^{-1}(X_s)$ sont affines
(\hyperref[II.1.2.5-fr]{1.2.5}) et recouvrent $X'$~; il résulte donc encore
du lemme (\hyperref[II.4.6.13.1-fr]{4.6.13.1, (i)}) et de
(\hyperref[II.4.5.2-fr]{4.5.2, a')}) que $\mathcal{L}'$ est ample, le
morphisme $f_{(Y')}$ étant quasi-compact
(\hyperref[I.6.6.4-fr]{I, 6.6.4, (iii)}).

Pour prouver (iv), notons d'abord que $f_1\times_S f_2$ est quasi-compact
(\hyperref[I.6.6.4-fr]{I, 6.6.4, (iv)}). On peut en outre supposer $S$,
$Y_1$ et $Y_2$ affines ((\hyperref[II.4.6.4-fr]{4.6.4}) et
(\hyperref[I.3.2.7-fr]{I, 3.2.7})), donc $\mathcal{L}_i$ ample ($i=1,2$)
(\hyperref[II.4.6.6-fr]{4.6.6}). Les ouverts
$(X_1)_{s_1}\times_S(X_2)_{s_2}$ forment un recouvrement de
$X_1\times_S X_2$ lorsque $s_i$ parcourt les sections des
$\mathcal{L}_i^{\otimes n_i}$ telles que $(X_i)_{s_i}$ soit affine
(\hyperref[II.4.5.2-fr]{4.5.2, a')}). D'ailleurs, en remplaçant $s_1$ et
$s_2$ par des puissances convenables, ce qui ne change pas les
$(X_i)_{s_i}$, on peut toujours supposer $n_1=n_2$. On déduit alors de
(\hyperref[II.4.6.13.1-fr]{4.6.13.1, (ii)}) et de
(\hyperref[II.4.5.2-fr]{4.5.2, a')}) que
$\mathcal{L}_1\otimes_S\mathcal{L}_2$ est ample, d'où l'assertion, puisque
$Y_1\times_S Y_2$ est affine (\hyperref[II.4.6.6-fr]{4.6.6}).

Reste à prouver (ii). Par le même raisonnement que dans
(\hyperref[II.4.4.10-fr]{4.4.10}), mais utilisant ici
(\hyperref[II.4.6.4-fr]{4.6.4}), on peut se borner au cas où $Z$ est
affine. Comme $\mathcal{K}$ est alors ample, et $Y$ quasi-compact, il existe
un nombre fini de sections $s_i\in\Gamma(Y,\mathcal{K}^{\otimes k_i})$
telles que les $Y_{s_i}$

